% Quadro gerado de research/data/mmat_reassessment_current.csv
% Reavaliação documental preliminar; não representa score nem avaliação final.
% Não editar manualmente; execute: python -m src.analysis.mmat_current_tcc_table
\begingroup
\SingleSpacing
\setlength{\parindent}{0pt}
\fontsize{10}{12}\selectfont
\def\LTcaptype{quadro}
\setlength{\tabcolsep}{2pt}
\setlength{\LTcapwidth}{\textwidth}
\begin{longtable}{|>{\raggedright\arraybackslash}p{4.5cm}|>{\raggedright\arraybackslash}p{4.5cm}|ccccccc|}
\caption{Triagem e apreciação preliminar pelo MMAT 2018 de 15 estudos do corpus.}\label{qua:mmat-reavaliacao-atual}\\
\hline
\textbf{Estudo} & \textbf{Desenho} & \textbf{S1} & \textbf{S2} & \textbf{Q1} & \textbf{Q2} & \textbf{Q3} & \textbf{Q4} & \textbf{Q5} \\
\hline
\endfirsthead
\multicolumn{9}{c}{\quadroname\ \thequadro\ -- Continuação}\\
\hline
\textbf{Estudo} & \textbf{Desenho} & \textbf{S1} & \textbf{S2} & \textbf{Q1} & \textbf{Q2} & \textbf{Q3} & \textbf{Q4} & \textbf{Q5} \\
\hline
\endhead
\hline
\multicolumn{9}{r}{\textit{Continua na próxima página}}\\
\endfoot
\hline
\endlastfoot
\citeonline{Math2021_001} & Quant. descritivo & Y & Y & Y & CT & CT & CT & CT \\
\citeonline{Implementation2025_000} & Quant. descritivo & Y & Y & Y & CT & CT & CT & Y \\
\citeonline{Multimodels2020_002} & Quant. descritivo & Y & Y & CT & CT & CT & CT & CT \\
\citeonline{Analysis2022_003} & Quant. descritivo & Y & Y & CT & CT & CT & CT & CT \\
\citeonline{Design2025_004} & Cat. 3 utilizada (alocação incerta) & Y & Y & CT & CT & CT & CT & CT \\
\citeonline{Identifying2017_006} & Quant. descritivo & Y & Y & Y & CT & CT & CT & CT \\
\citeonline{Innovative2023_005} & Cat. 3 utilizada (alocação incerta) & Y & Y & CT & CT & CT & CT & CT \\
\citeonline{Computational2017_008} & Cat. 4 adaptada (modelagem) & Y & Y & Y & CT & CT & CT & CT \\
\citeonline{Machine2024_009} & Quant. descritivo & Y & Y & CT & CT & CT & CT & CT \\
\citeonline{Enhancing2025_012} & Cat. 2 declarada (contraditória) & Y & Y & CT & N & Y & CT & CT \\
\citeonline{He2025_6915} & Quant. não rand. (componente escolar) & Y & Y & CT & Y & CT & N & Y \\
\citeonline{Villegas2025_6916} & Métodos mistos & Y & Y & Y & Y & CT & CT & CT \\
\citeonline{Ozseven2026_6917} & Quant. descritivo & Y & Y & Y & CT & CT & CT & Y \\
\citeonline{UniversityMathematics2026_6919} & Métodos mistos & Y & Y & Y & Y & N & CT & N \\
\citeonline{Echeveria2025_6920} & Métodos mistos & Y & Y & Y & Y & CT & CT & N \\
\end{longtable}
\fonte{Elaborado pelo autor (2026).}
\par\noindent\textit{Nota:} S1/S2 = perguntas de triagem; Q1--Q5 = critérios da categoria; Y = sim; N = não; CT = não é possível concluir; Quant. = quantitativo; rand. = randomizado. Ver Seção~\ref{sec:avaliacao-qualidade}.\par
\endgroup

% Quadro gerado de research/data/mmat_component_reassessment_current.csv
\begingroup
\SingleSpacing
\setlength{\parindent}{0pt}
\def\LTcaptype{quadro}
\fontsize{10}{12}\selectfont
\setlength{\tabcolsep}{2pt}
\renewcommand{\arraystretch}{1.15}
\begin{longtable}{|>{\raggedright\arraybackslash}p{2.2cm}|>{\raggedright\arraybackslash}p{2.2cm}|ccccc|>{\raggedright\arraybackslash}p{6.1cm}|}
\caption{Componentes qualitativos e quantitativos dos estudos de métodos mistos.}\label{qua:mmat-componentes}\\
\hline
\textbf{Estudo} & \textbf{Componente} & \textbf{C1} & \textbf{C2} & \textbf{C3} & \textbf{C4} & \textbf{C5} & \textbf{Síntese e páginas consultadas} \\
\hline
\endfirsthead
\multicolumn{8}{c}{\quadroname\ \thequadro\ -- Continuação}\\
\hline
\textbf{Estudo} & \textbf{Componente} & \textbf{C1} & \textbf{C2} & \textbf{C3} & \textbf{C4} & \textbf{C5} & \textbf{Síntese e páginas consultadas} \\
\hline
\endhead
\hline
\multicolumn{8}{r}{\textit{Continua na próxima página}}\\
\endfoot
\hline
\endlastfoot
\citeonline{Villegas2025_6916} & Quant. não rand. & CT & CT & CT & CT & CT & 3.1 (CT): A população de recrutamento e a estratégia de amostragem não estão estabelecidas nas seções examinadas. 3.2 (CT): As medidas de resultado são descritas, mas o relato disponível não fundamenta suficientemente sua adequação. 3.3 (CT): A completude dos dados de resultado e as perdas de participantes não estão estabelecidas. 3.4 (CT): A alocação dos grupos e o controle inicial de fatores de confusão não estão suficientemente documentados. 3.5 (CT): A introdução menciona um semestre acadêmico, os métodos relatam seis semanas e os resultados mostram progresso na oitava semana para estudantes de programação com alta interação; essa divergência não estabelece a adesão ou a cronologia do componente de matemática. Páginas e seções: Páginas físicas 2 (450 estudantes e semestre letivo), 13–16 (delineamento, medidas e análise), 18–23 (resultados; semana 8 na p. 22) e 24–27 (testes e limites) \\
\citeonline{Villegas2025_6916} & Qualitativo & Y & Y & CT & CT & CT & 1.1 (Y): A abordagem qualitativa é pertinente para examinar as percepções de estudantes e educadores sobre o sistema de tutoria. 1.2 (Y): Entrevistas estruturadas com estudantes e educadores, realizadas ao fim da experiência e com tópicos pertinentes explicitados, são adequadas para explorar as percepções sobre o sistema (página física 16). Detalhes de seleção e registro continuam como perguntas de esclarecimento, sem serem tratados automaticamente como inadequação da coleta. 1.3 (CT): O relato menciona análise temática e temas, mas não explicita suficientemente a derivação analítica a partir das entrevistas (páginas físicas 16–17). 1.4 (CT): As interpretações das entrevistas não são suficientemente vinculadas ao material qualitativo examinado (páginas físicas 16–17); médias e percentuais de satisfação nas pp. 24–25 são dados quantitativos e não resolvem esse vínculo. 1.5 (CT): A coerência entre fontes qualitativas, coleta, análise e interpretação permanece pouco clara. Páginas e seções: Páginas físicas 16–17 (Qualitative Analysis); 24–25 (satisfação quantitativa, não dados de entrevista) \\
\citeonline{UniversityMathematics2026_6919} & Quant. não rand. & CT & Y & CT & N & CT & 3.1 (CT): Os grupos pertencem a uma disciplina e estão separados por quatro anos. O artigo menciona amostragem por conglomerados, mas não define a população-alvo nem explica como as turmas de 2021 e 2025 foram escolhidas (páginas físicas 7–8). 3.2 (Y): O questionário e os resultados da disciplina são descritos e analisados; a apreciação permanece provisória. 3.3 (CT): A completude e o tratamento dos dados de resultado não estão suficientemente documentados nas seções examinadas. 3.4 (N): Os grupos de controle e experimental de anos diferentes introduzem confundimento temporal reconhecido pelos autores; os testes de linha de base não controlam mudanças de equipe, contexto e uso de IA (páginas físicas 8 e 24). 3.5 (CT): A implementação da intervenção e a adesão não estão suficientemente documentadas para este critério. Páginas e seções: Páginas físicas 7–8 (§4.1.1), 12–14 (§4.3–4.4.1, tabela 4) e 24 (§6) \\
\citeonline{UniversityMathematics2026_6919} & Qualitativo & Y & Y & CT & Y & CT & 1.1 (Y): As entrevistas são adequadas para explorar as perspectivas dos estudantes sobre a reforma da disciplina. 1.2 (Y): O artigo relata entrevistas semiestruturadas com 12 estudantes do grupo experimental, selecionados aleatoriamente logo após o curso (página física 18). 1.3 (CT): As etapas analíticas utilizadas para derivar os temas não estão suficientemente detalhadas. 1.4 (Y): A interpretação inclui relatos dos estudantes e é apoiada por citações apresentadas no artigo (página física 18). 1.5 (CT): A relação entre amostragem qualitativa, coleta, análise e conclusões integradas requer verificação adicional. Páginas e seções: Páginas físicas 18 (§4.4.8) e 24 (§6) \\
\citeonline{Echeveria2025_6920} & Quant. descritivo & Y & CT & N & CT & CT & 4.1 (Y): Os registros de estudantes de uma escola são adequados à questão de classificação local descrita no estudo. 4.2 (CT): A escola e seus estudantes delimitam uma população local; o MMAT não exige uma população-alvo mais abrangente. A correspondência dos 428 casos completos aos 633 recrutados não foi estabelecida, com anos letivos divergentes e sem comparação suficiente dos retidos e não retidos (páginas físicas 3 e 7). 4.3 (N): Oito itens da escala usados para construir as categorias de mentalidade também permanecem entre os preditores selecionados, comprometendo a mensuração independente do alvo. 4.4 (CT): O relato informa 633 participantes recrutados e 428 casos completos e cita os anos letivos 2023–2024 e 2022–2023 sem esclarecer a origem dos grupos nem se os casos ausentes produzem viés de não resposta. 4.5 (CT): As descrições da divisão dos dados e da validação divergem e não permitem apreciar claramente a adequação da análise. Páginas e seções: Páginas físicas 3/\allowbreak{}impressa 65 (§2.1–2.3), 6/\allowbreak{}68 (§2.6) e 7–8/\allowbreak{}69–70 (resultados) \\
\citeonline{Echeveria2025_6920} & Qualitativo & Y & Y & CT & Y & CT & 1.1 (Y): As entrevistas com professores são pertinentes aos objetivos declarados de contextualização e triangulação do estudo. 1.2 (Y): O relato identifica entrevistas com cinco docentes de Matemática para explorar o uso do modelo e apresenta suas falas individuais (páginas físicas 3 e 11–12), constituindo coleta pertinente à questão de percepção docente. A seleção e o registro das respostas permanecem pouco detalhados. 1.3 (CT): O procedimento analítico do material das entrevistas não está adequadamente descrito. 1.4 (Y): As falas individuais dos docentes sustentam a interpretação descritiva de suas percepções sobre identificação, apoio e planejamento (páginas físicas 11–12). O julgamento não transforma expectativas dos entrevistados em efeito observado sobre aprendizagem. 1.5 (CT): A coerência entre fontes, coleta, análise e interpretação qualitativas não pode ser concluída, pois o procedimento de análise das entrevistas permanece pouco explicitado (páginas físicas 3 e 11–13). A integração com dados quantitativos é examinada separadamente nos critérios de métodos mistos. Páginas e seções: Páginas físicas 3/\allowbreak{}impressa 65 (§2.1 e coleta), 11–12/\allowbreak{}73–74 (falas docentes) e 13/\allowbreak{}75 (conclusões) \\
\end{longtable}
\fonte{Elaborado pelo autor (2026).}
\par\noindent\textit{Nota:} C1--C5 = critérios da categoria; Y = sim; N = não; CT = não é possível concluir; Quant. = quantitativo; rand. = randomizado. Ver Seção~\ref{sec:avaliacao-qualidade}.\par
\endgroup
